% Zone „Das kennst du schon“ – aus bank/pyramide-kegel-kugel/zone.jsonl (zusammenbau v0.1)
\einheitenkopf[zone][Kennst du schon]{Das kennst du schon}
\setcounter{aufgabe}{0}

%% TODO Titel: Ich-kann-Satz fehlt in der Bank (Platzhalter: Kettenname) – Nr. 1
%% TODO Anweisung: ein Satz über der ersten Teilaufgabe fehlt in der Bank – Nr. 1
\begin{aufgabe}{Flächeninhalt von Quadrat, Rechteck, Dreieck (Seitenflächen) und Kreis (Grundkreis)}
\begin{teile}
\teil Quadrat mit der Seite $7$ cm – Flächeninhalt? \leerfeld[cm²]
\teil Kreis mit dem Radius $3$ cm – Flächeninhalt? Runde auf eine Stelle. \leerfeld[cm²]
\end{teile}
\end{aufgabe}

%% TODO Titel: Ich-kann-Satz fehlt in der Bank (Platzhalter: Kettenname) – Nr. 2
%% TODO Anweisung: ein Satz über der ersten Teilaufgabe fehlt in der Bank – Nr. 2
\begin{aufgabe}{Volumen von Prisma und Zylinder als Grundfläche mal Höhe (der Bezugskörper für das Drittel) und Mantel des Zylinders}
\begin{teile}
\teil Prisma mit der Grundfläche $8$ cm² und der Höhe $6$ cm – Volumen? \leerfeld[cm³]
\teil Zylinder mit dem Radius $2$ cm und der Höhe $9$ cm – Volumen? Runde auf eine Stelle. \leerfeld[cm³]
\end{teile}
\end{aufgabe}

%% TODO Titel: Ich-kann-Satz fehlt in der Bank (Platzhalter: Kettenname) – Nr. 3
%% TODO Anweisung: ein Satz über der ersten Teilaufgabe fehlt in der Bank – Nr. 3
\begin{aufgabe}{Satz des Pythagoras}
\begin{teile}
\teil Rechtwinkliges Dreieck, Katheten $3$ cm und $4$ cm – Hypotenuse? \leerfeld[cm]
\teil Rechtwinkliges Dreieck, Katheten $1{,}2$ m und $0{,}9$ m – Hypotenuse? \leerfeld[m]
\end{teile}
\end{aufgabe}

%% TODO Titel: Ich-kann-Satz fehlt in der Bank (Platzhalter: Kettenname) – Nr. 4
%% TODO Anweisung: ein Satz über der ersten Teilaufgabe fehlt in der Bank – Nr. 4
\begin{aufgabe}{Quadrieren und hoch drei, Wurzel und Kubikwurzel mit dem Taschenrechner, Rechnen mit π, Runden auf eine sinnvolle Stelle}
\begin{teile}
\teil $3^3$ – wie viel? \leerfeld
\teil $1{,}5^3$ – wie viel? \leerfeld
\end{teile}
\end{aufgabe}

%% TODO Titel: Ich-kann-Satz fehlt in der Bank (Platzhalter: Kettenname) – Nr. 5
%% TODO Anweisung: ein Satz über der ersten Teilaufgabe fehlt in der Bank – Nr. 5
\begin{aufgabe}{Ein Drittel und vier Drittel einer Zahl}
\begin{teile}
\teil Ein Drittel von $27$? \leerfeld
\teil Vier Drittel von $4{,}5$? \leerfeld
\end{teile}
\end{aufgabe}

%% TODO Titel: Ich-kann-Satz fehlt in der Bank (Platzhalter: Kettenname) – Nr. 6
%% TODO Anweisung: ein Satz über der ersten Teilaufgabe fehlt in der Bank – Nr. 6
\begin{aufgabe}{Formel aus der Formelsammlung entnehmen und nach einer Größe umstellen (Volumenformel mit Bruchfaktor nach der Höhe)}
\begin{teile}
\teil Rechteck: $A = a \cdot b$ mit $A = 42$ m² und $a = 6$ m. Wie lang ist $b$? $b$ = \leerfeld[m]
\teil Prisma: $V = G \cdot h$ mit dem Volumen $V = 37{,}5$ dm³ und der Grundfläche $G = 7{,}5$ dm². Wie hoch ist das Prisma? \leerfeld[dm]
\end{teile}
\end{aufgabe}

%% TODO Titel: Ich-kann-Satz fehlt in der Bank (Platzhalter: Kettenname) – Nr. 7
%% TODO Anweisung: ein Satz über der ersten Teilaufgabe fehlt in der Bank – Nr. 7
\begin{aufgabe}{Radius aus Durchmesser}
\begin{teile}
\teil Durchmesser $18$ cm – Radius? \leerfeld[cm]
\teil Durchmesser $45$ mm – Radius in cm? \leerfeld[cm]
\end{teile}
\end{aufgabe}

%% TODO Titel: Ich-kann-Satz fehlt in der Bank (Platzhalter: Kettenname) – Nr. 8
%% TODO Anweisung: ein Satz über der ersten Teilaufgabe fehlt in der Bank – Nr. 8
\begin{aufgabe}{Volumeneinheiten cm³ → dm³ = l mit 1000, m³; Masse aus Volumen und Dichte}
\begin{teile}
\teil $3\,000$ cm³ – wie viel Liter? \leerfeld[l]
\teil $750$ cm³ – wie viel Liter? \leerfeld[l]
\end{teile}
\end{aufgabe}

%% TODO Titel: Ich-kann-Satz fehlt in der Bank (Platzhalter: Kettenname) – Nr. 9
%% TODO Anweisung: ein Satz über der ersten Teilaufgabe fehlt in der Bank – Nr. 9
\begin{aufgabe}{Termumformung mit Variablen}
\begin{teile}
\teil $(2a)^2$ – vereinfache. \leerfeld
\teil $(0{,}5b)^2$ – vereinfache. \leerfeld
\end{teile}
\end{aufgabe}

%% TODO Titel: Ich-kann-Satz fehlt in der Bank (Platzhalter: Kettenname) – Nr. 10
%% TODO Anweisung: ein Satz über der ersten Teilaufgabe fehlt in der Bank – Nr. 10
\begin{aufgabe}{Radius aus Durchmesser – Fehler finden}
\begin{teile}
\teil Kreis mit dem Durchmesser $16$ cm. Lena rechnet: \rechnung{A &= \pi \cdot 16^2 \\ A &\approx 804{,}2 \text{ cm}^2} Finde den Fehler.
\end{teile}
\end{aufgabe}

%% TODO Titel: Ich-kann-Satz fehlt in der Bank (Platzhalter: Kettenname) – Nr. 11
%% TODO Anweisung: ein Satz über der ersten Teilaufgabe fehlt in der Bank – Nr. 11
\begin{aufgabe}{Radius aus Durchmesser – selbst rechnen}
\begin{teile}
\teil Kreis mit dem Durchmesser $14$ cm – Flächeninhalt? Runde auf eine Stelle. \leerfeld[cm²]
\end{teile}
\end{aufgabe}
